\section{Document Quality — LLM-as-a-Judge}
\label{sec:eval-judge}

The LLM-as-a-Judge evaluation provides a reference-free, GDPR-specific quality
signal across five dimensions: \emph{Completeness}, \emph{Scenario
Faithfulness}, \emph{Legal Correctness}, \emph{Hallucination} (inverse-coded;
higher equals less hallucination), and \emph{Overall}. The judge is a single
LLM (Gemini~3.1 Pro Preview, temperature $0$) prompted with the rubric
documented in Appendix~\ref{app:judge-prompt}. All $113$ retained documents
are scored. Per-mode results are reported in the order Form, Ask, Chat,
followed by the aggregated mean across all documents.

\subsection{Per-Dimension Results}

Table~\ref{tab:judge-summary} reports the per-dimension mean and standard
deviation per mode and the aggregated mean. Figure~\ref{fig:judge-boxplot}
shows the corresponding distributions.

\begin{table}[htbp]
\centering
\caption{LLM-as-Judge dimension means $\pm$ standard deviation by mode (full
sample $n = 113$, scale $1$--$5$, higher is better including for
\emph{Hallucination (inverse)}). The aggregated ``All'' column is the mean
across all $113$ documents.}
\label{tab:judge-summary}
\begin{tabular}{@{}l r r r r@{}}
\toprule
Dimension & Form & Ask & Chat & All \\
\midrule
Completeness            & $4.84 \pm 0.55$ & $3.84 \pm 1.54$ & $3.18 \pm 1.59$ & $3.94 \pm 1.43$ \\
Scenario Faithfulness   & $1.84 \pm 0.76$ & $2.30 \pm 1.20$ & $2.49 \pm 1.19$ & $2.21 \pm 1.10$ \\
Legal Correctness       & $3.59 \pm 0.55$ & $3.43 \pm 1.34$ & $3.15 \pm 1.42$ & $3.39 \pm 1.18$ \\
Hallucination (inverse) & $1.65 \pm 0.72$ & $2.40 \pm 1.12$ & $2.54 \pm 0.91$ & $2.20 \pm 1.02$ \\
Overall                 & $2.41 \pm 0.69$ & $2.49 \pm 1.04$ & $2.62 \pm 1.31$ & $2.50 \pm 1.04$ \\
\bottomrule
\end{tabular}
\end{table}

\begin{figure}[htbp]
\centering
\includegraphics[width=\linewidth]{judge_metrics_boxplot.png}
\caption{Distributions of the five LLM-as-Judge dimensions by mode (full
sample $n = 113$, scale $1$--$5$).}
\label{fig:judge-boxplot}
\end{figure}

For \emph{Completeness}, Form scores highest by a substantial margin (mean
$4.84$, median $5.0$), with Ask second ($3.84$, median $5.0$) and Chat lowest
($3.18$, median $3.0$); aggregated mean $3.94$. Completeness decreases by
approximately $1.66$~points from Form to Chat. Form's standard deviation
is also the lowest of the three modes ($0.55$ versus Ask $1.54$ and Chat
$1.59$), indicating that Form yields uniformly complete documents.

For \emph{Scenario Faithfulness}, the ordering inverts. Form scores lowest
($1.84$, median $2.0$), Ask second ($2.30$, median $2.0$), Chat highest
($2.49$, median $2.0$); aggregated mean $2.21$. Faithfulness to the specific
study scenario increases by approximately $0.65$~points from Form to Chat.

For \emph{Legal Correctness}, Form scores highest ($3.59$, median $4.0$),
Ask second ($3.43$, median $4.0$), Chat lowest ($3.15$, median $3.0$);
aggregated mean $3.39$. The spread is approximately $0.45$~points across
modes; Form's standard deviation ($0.55$) is again the lowest, reflecting a
narrow concentration of Legal Correctness scores around the value $4.0$.

For \emph{Hallucination (inverse)}, the ordering matches Scenario
Faithfulness. Form scores lowest ($1.65$, median $2.0$), Ask second ($2.40$,
median $2.0$), Chat highest ($2.54$, median $3.0$); aggregated mean $2.20$.
Because the dimension is inverse-coded, this means Form documents are judged
to contain the most fabricated content and Chat documents the least.

For \emph{Overall}, the three modes are clustered closely together: Form
$2.41$, Ask $2.49$, Chat $2.62$; aggregated $2.50$. The aggregate Overall
score thus sits below the midpoint of the $1$--$5$ scale across all modes.
Form's standard deviation ($0.69$) is again the lowest, while Chat's is the
highest ($1.31$), indicating that Chat produces both the highest- and
lowest-rated documents in the sample.

\subsection{Cross-Dimension Patterns}

Three patterns emerge from the per-dimension means. First, the
Completeness ordering (Form $>$ Ask $>$ Chat) is the inverse of the
Scenario Faithfulness and Hallucination orderings (Form $<$ Ask $<$ Chat).
Form documents are judged to be the most structurally complete but also
the most prone to fabricated content, while Chat documents are judged to be
the most faithful to the specific scenario but the least structurally
complete. Ask sits between the two extremes on every dimension.

Second, the Legal Correctness scores are flat relative to the Completeness
spread: the gap between modes is approximately $0.45$~points on Legal
Correctness compared to $1.66$~points on Completeness. The judge therefore
distinguishes between modes more sharply on whether the document covers
all required Article 30(1) fields than on whether the legal terminology is
appropriate.

Third, the Overall scores compress the per-dimension spread into a narrow
range (Form $2.41$ to Chat $2.62$, $\Delta = 0.21$), with the ordering Form
$<$ Ask $<$ Chat. This is consistent with Overall reflecting a holistic
judgement that weighs Faithfulness and Hallucination against Completeness;
Form's high Completeness is partially offset by its low Faithfulness and
high Hallucination, and the net result is a slight Chat advantage.

\subsection{Note on the Form Completeness Pattern}

A descriptive observation that recurs across the judge data is that
$31$ of the $37$ Form documents ($83.8\,\%$) achieve Completeness $= 5$ while
simultaneously scoring Hallucination (inverse) $\leq 2$, indicating
``moderate'' or ``significant'' hallucination per the judge rubric. Of these,
$17$ Form documents combine the maximum Completeness score of $5$ with the
minimum Hallucination (inverse) score of $1$. By contrast, the Ask and Chat
distributions on Completeness are more dispersed (standard deviations of
$1.54$ and $1.59$ respectively versus Form's $0.55$), and the joint
incidence of high Completeness with low Hallucination (inverse) is
correspondingly less frequent. This descriptive co-occurrence is the
proximate source of the Form-versus-Faithfulness inversion noted above; its
interpretation in terms of the checkbox interface design is reserved for
the Discussion chapter.
